\documentclass[letterpaper,12pt]{article}

\usepackage{fancyhdr}
\usepackage[most]{tcolorbox}
\usepackage{longtable}
\usepackage{makecell}
\usepackage{tikz}
\usepackage{geometry}
\usepackage{array}
\usepackage{multicol}
\usepackage[table]{xcolor}
\usepackage{polyglossia}
\setmainlanguage{english}
\setotherlanguage{arabic}
\setotherlanguage{hebrew}
\usepackage{fontspec}
\setmainfont{Charis SIL}
\newfontfamily\arabicfont[Script=Arabic,Scale=1.2]{Amiri}
\newfontfamily\hebrewfont{Ezra SIL}[Script=Hebrew]
\geometry{margin=1.5cm}
\usepackage{booktabs}
\usepackage{graphicx}
\usepackage{multirow}

\setlength{\headheight}{14.49998pt}
\addtolength{\topmargin}{-2.49998pt}
\pagestyle{fancy}
\fancyhf{}
\fancyfoot[C]{\thepage}
\fancyhead[R]{\textit{Why Should I Learn Arabic?}}
\renewcommand{\headrulewidth}{0pt}

\usetikzlibrary{shapes,arrows,decorations.pathmorphing}

% Custom colors
\definecolor{headercolor}{RGB}{70,130,180}
\definecolor{boxcolor}{RGB}{240,248,255}
\definecolor{accentcolor}{RGB}{220,20,60}
\definecolor{tableheader}{RGB}{220,220,220}
\definecolor{dialectcolor}{RGB}{34,139,34}

\begin{document}

\title{\textbf{\Large Arabic Phrase Analyser}\\
\large Why Should I Learn This Arabic Language?\\
\normalsize \textarabic{لم يجب علي تعلم ذي اللغة العربية}}
\author{Igor Deruga}
\date{}
\maketitle

% ======================== Phrase Display ========================
\begin{tcolorbox}[colback=boxcolor,colframe=headercolor,title=\textbf{Arabic Phrase with Full Diacritics},breakable]
\centering
\Large\textarabic{لِمَ يَجِبُ عَلَيَّ تَعَلُّمُ ذِي اللُّغَةِ الْعَرَبِيَّةِ؟}
\end{tcolorbox}

% ======================== Translations ========================
\section{English Translation}
\begin{tcolorbox}[colback=white,colframe=accentcolor,breakable]
\textbf{Literal:} For what must upon-me learning of this the-language the-Arabic? \\
\textit{[Arabic order retained for direct mapping]}\\[0.5em]
\textbf{Adapted:} Why should I learn this Arabic language? / Why must I learn Arabic?
\end{tcolorbox}

% ======================== Detailed Word Analysis ========================
\section{Detailed Word Analysis}

\subsection{\textarabic{لِمَ} – [lima]}
\begin{tabular}{p{3cm}p{10cm}}
\toprule
\textbf{Translation} & why / for what \\
\textbf{Root} & \textarabic{م-ا} (question particle) \\
\textbf{Pattern} & Interrogative particle \textarabic{مَا} with preposition \textarabic{لِ} \\
\textbf{Grammar} & Interrogative compound: \textarabic{لِ} (for) + \textarabic{مَا} (what) with elision of alif \\
\midrule
\textbf{Formation} & \textarabic{لِمَ} →  \textarabic{ل} + \textarabic{مَا} (when \textarabic{مَا} follows a preposition, the alif drops) \\
\midrule
\textbf{Examples} & \makecell[l]{\parbox{9.5cm}{
1. \textarabic{لِمَ تَبْكِي؟} - Why do you cry? [lima tabkiː]\\
2. \textarabic{لِمَ جِئْتَ مُبَكِّرًا؟} - Why did you come early? [lima ʤiʔta muˈbakkiran]\\
3. \textarabic{لِمَ لَا تَدْرُسُ؟} - Why don't you study? [lima laː tadrusu]
}} \\
\midrule
\textbf{Synonyms} & \textarabic{لِمَاذَا} (limāðā), \textarabic{لِأَيِّ شَيْءٍ} (for which thing) \\
\textbf{Etymology} & Combination of preposition \textarabic{لِ} (for/to) + interrogative \textarabic{مَا} (what) \\
\bottomrule
\end{tabular}

\subsubsection*{Usage Notes}
\begin{itemize}
  \item \textarabic{لِمَ} is more literary/formal than \textarabic{لِمَاذَا}
  \item The alif of \textarabic{مَا} elides after prepositions: \textarabic{لِمَ، فِيمَ، عَلَامَ، إِلَامَ}
  \item Can be followed by verbs in different tenses depending on context
  \item Often implies seeking justification or reason
\end{itemize}

\subsection{\textarabic{يَجِبُ} – [jaʤibu]}

\begin{tabular}{p{3cm}p{10cm}}
\toprule
\textbf{Translation} & must / is necessary / should \\
\textbf{Root} & \textarabic{و-ج-ب} (w-ʤ-b) \\
\textbf{Pattern} & \textarabic{يَفْعِلُ} (jafʕilu) - Form I imperfect \\
\textbf{Grammar} & Present tense verb, 3rd person masculine singular, indicative mood \\
\midrule
\textbf{Examples} & \makecell[l]{\parbox{9.5cm}{
1. \textarabic{يَجِبُ عَلَيْكَ أَنْ تَذْهَبَ} - You must go [jaʤibu ʕalajka ʔan taðhaba]\\
2. \textarabic{وَجَبَ الصَّلَاةُ} - Prayer became obligatory [waʤaba sˤ-sˤalaːtu]\\
3. \textarabic{سَيَجِبُ عَلَيْهِ الدَّفْعُ} - He will have to pay [sajaʤibu ʕalajhi d-dafʕu]
}} \\
\midrule
\textbf{Synonyms} & \textarabic{يَلْزَمُ} (is binding), \textarabic{يَنْبَغِي} (should), \textarabic{عَلَى} (upon/must) \\
\textbf{Etymology} & From root meaning "to be necessary, to be incumbent, to fall due" \\
\bottomrule
\end{tabular}

\subsection{Conjugation}
\par{\large \textbf{Verbal Noun}: \textarabic{وُجُوب} /wuʤuːb/ - necessity, obligation}

\begin{longtable}{|>{\raggedright}p{3.5cm}|p{5cm}|p{5cm}|}
\hline
\textbf{Person} & \textbf{Perfect (Past)} & \textbf{Imperfect (Present)} \\
\hline
\textbf{3rd person masculine singular} & \textarabic{وَجَبَ} /waʤaba/ & \textarabic{يَجِبُ} /jaʤibu/ \\
\hline
\textbf{3rd person feminine singular} & \textarabic{وَجَبَتْ} /waʤabat/ & \textarabic{تَجِبُ} /taʤibu/ \\
\hline
\textbf{3rd person masculine dual} & \textarabic{وَجَبَا} /waʤabaː/ & \textarabic{يَجِبَانِ} /jaʤibaːni/ \\
\hline
\textbf{3rd person feminine dual} & \textarabic{وَجَبَتَا} /waʤabataː/ & \textarabic{تَجِبَانِ} /taʤibaːni/ \\
\hline
\textbf{3rd person masculine plural} & \textarabic{وَجَبُوا} /waʤabuː/ & \textarabic{يَجِبُونَ} /jaʤibuːna/ \\
\hline
\textbf{3rd person feminine plural} & \textarabic{وَجَبْنَ} /waʤabna/ & \textarabic{يَجِبْنَ} /jaʤibna/ \\
\hline
\textbf{2nd person masculine singular} & \textarabic{وَجَبْتَ} /waʤabta/ & \textarabic{تَجِبُ} /taʤibu/ \\
\hline
\textbf{2nd person feminine singular} & \textarabic{وَجَبْتِ} /waʤabti/ & \textarabic{تَجِبِينَ} /taʤibiːna/ \\
\hline
\textbf{2nd person dual (m./f.)} & \textarabic{وَجَبْتُمَا} /waʤabtumaː/ & \textarabic{تَجِبَانِ} /taʤibaːni/ \\
\hline
\textbf{2nd person masculine plural} & \textarabic{وَجَبْتُمْ} /waʤabtum/ & \textarabic{تَجِبُونَ} /taʤibuːna/ \\
\hline
\textbf{2nd person feminine plural} & \textarabic{وَجَبْتُنَّ} /waʤabtunna/ & \textarabic{تَجِبْنَ} /taʤibna/ \\
\hline
\textbf{1st person singular} & \textarabic{وَجَبْتُ} /waʤabtu/ & \textarabic{أَجِبُ} /ʔaʤibu/ \\
\hline
\textbf{1st person plural} & \textarabic{وَجَبْنَا} /waʤabnaː/ & \textarabic{نَجِبُ} /naʤibu/ \\
\hline
\end{longtable}

\subsubsection*{Conjugation Notes}
\begin{itemize}
  \item \textbf{Impersonal usage}: \textarabic{يَجِبُ} is commonly used impersonally with \textarabic{عَلَى} + pronoun suffix to express obligation
  \item \textbf{With infinitive}: \textarabic{يَجِبُ أَنْ + subjunctive} - "must/should (do)"
  \item \textbf{Passive voice}: 
    \begin{itemize}
      \item Perfect: \textarabic{وُجِبَ} [wuʤiba] – it was obligated
      \item Imperfect: \textarabic{يُوجَبُ} [juːʤabu] – it is obligated
    \end{itemize}
  \item \textbf{Related forms}:
    \begin{itemize}
      \item Form IV: \textarabic{أَوْجَبَ} [ʔawʤaba] – to necessitate, make obligatory
      \item Form X: \textarabic{اِسْتَوْجَبَ} [istawʤaba] – to deserve, merit
    \end{itemize}
\end{itemize}

\subsection{\textarabic{عَلَيَّ} – [ʕalajja]}

\begin{tabular}{p{3cm}p{10cm}}
\toprule
\textbf{Translation} & upon me / on me \\
\textbf{Root} & \textarabic{ع-ل-و} (ʕ-l-w) \\
\textbf{Pattern} & Preposition \textarabic{عَلَى} + 1st person singular pronoun suffix \\
\textbf{Grammar} & Prepositional phrase with attached pronoun, expresses obligation \\
\midrule
\textbf{Formation} & \textarabic{عَلَى + ي → عَلَيَّ} (the yāʔ is doubled: \textarabic{ـِيَّ}) \\
\midrule
\textbf{Examples} & \makecell[l]{\parbox{9.5cm}{
1. \textarabic{عَلَيَّ دَيْنٌ كَبِيرٌ} - I have a big debt (upon me) [ʕalajja dajnun kabiːrun]\\
2. \textarabic{يَجِبُ عَلَيَّ أَنْ أَعْمَلَ} - I must work [jaʤibu ʕalajja ʔan ʔaʕmala]\\
3. \textarabic{عَلَيَّ مَسْؤُولِيَّةٌ} - I have responsibility [ʕalajja masʔuːlijjatun]
}} \\
\midrule
\textbf{Synonyms} & \textarabic{يَنْبَغِي لِي} (I should), \textarabic{يَلْزَمُنِي} (it is binding on me) \\
\textbf{Etymology} & From root meaning "to be high, elevated"; preposition indicates burden/obligation \\
\bottomrule
\end{tabular}

\subsubsection*{Full Pronominal Declension of \textarabic{عَلَى}}

\begin{tabular}{|c|c|c|}
\hline
\textbf{Person} & \textbf{Singular} & \textbf{Plural} \\
\hline
1st person & \textarabic{عَلَيَّ} /ʕalajja/ & \textarabic{عَلَيْنَا} /ʕalajnaː/ \\
\hline
2nd person masc. & \textarabic{عَلَيْكَ} /ʕalajka/ & \textarabic{عَلَيْكُمْ} /ʕalajkum/ \\
\hline
2nd person fem. & \textarabic{عَلَيْكِ} /ʕalajki/ & \textarabic{عَلَيْكُنَّ} /ʕalajkunna/ \\
\hline
3rd person masc. & \textarabic{عَلَيْهِ} /ʕalajhi/ & \textarabic{عَلَيْهِمْ} /ʕalajhim/ \\
\hline
3rd person fem. & \textarabic{عَلَيْهَا} /ʕalajhaː/ & \textarabic{عَلَيْهِنَّ} /ʕalajhinna/ \\
\hline
\multicolumn{3}{|c|}{\textbf{Dual Forms}} \\
\hline
2nd person dual & \multicolumn{2}{c|}{\textarabic{عَلَيْكُمَا} /ʕalajkumaː/} \\
\hline
3rd person masc. dual & \multicolumn{2}{c|}{\textarabic{عَلَيْهِمَا} /ʕalajhimaː/} \\
\hline
3rd person fem. dual & \multicolumn{2}{c|}{\textarabic{عَلَيْهِمَا} /ʕalajhimaː/} \\
\hline
\end{tabular}

\subsection{\textarabic{تَعَلُّمُ} – [taʕallumu]}

\begin{tabular}{p{3cm}p{10cm}}
\toprule
\textbf{Translation} & learning / studying \\
\textbf{Root} & \textarabic{ع-ل-م} (ʕ-l-m) \\
\textbf{Pattern} & \textarabic{تَفَعُّلٌ} (tafaʕʕulun) - Form V verbal noun \\
\textbf{Grammar} & Verbal noun (maṣdar), nominative case, functioning as subject of \textarabic{يَجِبُ} \\
\midrule
\textbf{Form V Meaning} & Reflexive/intensive: to learn (for oneself), to study \\
\midrule
\textbf{Examples} & \makecell[l]{\parbox{9.5cm}{
1. \textarabic{تَعَلُّمُ اللُّغَاتِ مُفِيدٌ} - Learning languages is useful [taʕallumu l-luɣaːti mufiːdun]\\
2. \textarabic{أَحِبُّ تَعَلُّمَ الْعِلْمِ} - I love learning knowledge [ʔuħibbu taʕalluma l-ʕilmi]\\
3. \textarabic{تَعَلُّمُهُ سَهْلٌ} - Learning it is easy [taʕallumuhu sahlun]
}} \\
\midrule
\textbf{Synonyms} & \textarabic{دِرَاسَة} (study), \textarabic{تَحْصِيل} (acquisition), \textarabic{اِكْتِسَاب} (gaining) \\
\textbf{Etymology} & From root meaning "to know, to mark, to teach" (related to Hebrew \texthebrew{למד} /lamad/) \\
\bottomrule
\end{tabular}

\subsection{Verb Conjugation: \textarabic{تَعَلَّمَ} (Form V)}
\par{\large \textbf{Verbal Noun}: \textarabic{تَعَلُّم} /taʕallum/ - learning}

\begin{longtable}{|>{\raggedright}p{3.5cm}|p{5cm}|p{5cm}|}
\hline
\textbf{Person} & \textbf{Perfect (Past)} & \textbf{Imperfect (Present)} \\
\hline
\textbf{3rd person masculine singular} & \textarabic{تَعَلَّمَ} /taʕallama/ & \textarabic{يَتَعَلَّمُ} /jataʕallamu/ \\
\hline
\textbf{3rd person feminine singular} & \textarabic{تَعَلَّمَتْ} /taʕallamat/ & \textarabic{تَتَعَلَّمُ} /tataʕallamu/ \\
\hline
\textbf{3rd person masculine dual} & \textarabic{تَعَلَّمَا} /taʕallamaː/ & \textarabic{يَتَعَلَّمَانِ} /jataʕallamaːni/ \\
\hline
\textbf{3rd person feminine dual} & \textarabic{تَعَلَّمَتَا} /taʕallamataː/ & \textarabic{تَتَعَلَّمَانِ} /tataʕallamaːni/ \\
\hline
\textbf{3rd person masculine plural} & \textarabic{تَعَلَّمُوا} /taʕallamuː/ & \textarabic{يَتَعَلَّمُونَ} /jataʕallamuːna/ \\
\hline
\textbf{3rd person feminine plural} & \textarabic{تَعَلَّمْنَ} /taʕallamna/ & \textarabic{يَتَعَلَّمْنَ} /jataʕallamna/ \\
\hline
\textbf{2nd person masculine singular} & \textarabic{تَعَلَّمْتَ} /taʕallamta/ & \textarabic{تَتَعَلَّمُ} /tataʕallamu/ \\
\hline
\textbf{2nd person feminine singular} & \textarabic{تَعَلَّمْتِ} /taʕallamti/ & \textarabic{تَتَعَلَّمِينَ} /tataʕallamiːna/ \\
\hline
\textbf{2nd person dual (m./f.)} & \textarabic{تَعَلَّمْتُمَا} /taʕallamtumaː/ & \textarabic{تَتَعَلَّمَانِ} /tataʕallamaːni/ \\
\hline
\textbf{2nd person masculine plural} & \textarabic{تَعَلَّمْتُمْ} /taʕallamtum/ & \textarabic{تَتَعَلَّمُونَ} /tataʕallamuːna/ \\
\hline
\textbf{2nd person feminine plural} & \textarabic{تَعَلَّمْتُنَّ} /taʕallamtunna/ & \textarabic{تَتَعَلَّمْنَ} /tataʕallamna/ \\
\hline
\textbf{1st person singular} & \textarabic{تَعَلَّمْتُ} /taʕallamtu/ & \textarabic{أَتَعَلَّمُ} /ʔataʕallamu/ \\
\hline
\textbf{1st person plural} & \textarabic{تَعَلَّمْنَا} /taʕallamnaː/ & \textarabic{نَتَعَلَّمُ} /nataʕallamu/ \\
\hline
\end{longtable}

\subsubsection*{Related Forms from Root \textarabic{ع-ل-م}}
\begin{itemize}
  \item Form I: \textarabic{عَلِمَ} [ʕalima] – to know
  \item Form II: \textarabic{عَلَّمَ} [ʕallama] – to teach
  \item Form IV: \textarabic{أَعْلَمَ} [ʔaʕlama] – to inform, notify
  \item Form V: \textarabic{تَعَلَّمَ} [taʕallama] – to learn (reflexive)
  \item Form VI: \textarabic{تَعَالَمَ} [taʕaːlama] – to feign knowledge
  \item Form VIII: \textarabic{اِعْتَلَمَ} [iʕtalama] – rare variant of "to know"
\end{itemize}

\subsection{\textarabic{ذِي} – [ðiː]}

\begin{tabular}{p{3cm}p{10cm}}
\toprule
\textbf{Translation} & this (feminine, non-standard) \\
\textbf{Root} & Demonstrative pronoun \\
\textbf{Pattern} & \textarabic{ذِي} (non-standard for \textarabic{هٰذِهِ}) \\
\textbf{Grammar} & Demonstrative pronoun, genitive case (after \textarabic{تَعَلُّمُ}) \\
\midrule
\textbf{Note} & \makecell[l]{\parbox{9.5cm}{
\textbf{Important}: \textarabic{ذِي} is \textbf{dialectal/non-standard}. Classical Arabic uses:\\
• \textarabic{هٰذِهِ} [haːðihi] for feminine singular\\
The phrase should be: \textarabic{لِمَ يَجِبُ عَلَيَّ تَعَلُّمُ هٰذِهِ اللُّغَةِ الْعَرَبِيَّةِ؟}
}} \\
\midrule
\textbf{Examples} & \makecell[l]{\parbox{9.5cm}{
\textbf{Standard forms:}\\
1. \textarabic{هٰذَا كِتَابٌ} - This is a book (masc.) [haːðaː kitaːbun]\\
2. \textarabic{هٰذِهِ مَدْرَسَةٌ} - This is a school (fem.) [haːðihi madrasatun]\\
3. \textarabic{هٰؤُلَاءِ طُلَّابٌ} - These are students [haːʔulaːʔi tˤullaːbun]
}} \\
\midrule
\textbf{Dialectal Usage} & In some dialects, \textarabic{ذِي/دِي} replaces \textarabic{هٰذِهِ} \\
\textbf{Etymology} & From demonstrative root; related to \textarabic{ذَا، ذِهِ، ذَلِكَ} \\
\bottomrule
\end{tabular}

\subsubsection*{Standard Demonstrative Pronouns}

\begin{tabular}{|c|c|c|c|}
\hline
\textbf{Number} & \textbf{Masculine} & \textbf{Feminine} & \textbf{Common Plural} \\
\hline
Near (this/these) & \textarabic{هٰذَا} /haːðaː/ & \textarabic{هٰذِهِ} /haːðihi/ & \textarabic{هٰؤُلَاءِ} /haːʔulaːʔi/ \\
\hline
Far (that/those) & \textarabic{ذٰلِكَ} /ðaːlika/ & \textarabic{تِلْكَ} /tilka/ & \textarabic{أُولٰئِكَ} /ʔulaːʔika/ \\
\hline
Dual (these two) & \textarabic{هٰذَانِ} /haːðaːni/ & \textarabic{هَاتَانِ} /haːtaːni/ & — \\
\hline
Dual (those two) & \textarabic{ذَانِكَ} /ðaːnika/ & \textarabic{تَانِكَ} /taːnika/ & — \\
\hline
\end{tabular}

\subsection{\textarabic{اللُّغَةِ} – [al-luɣati]}

\begin{tabular}{p{3cm}p{10cm}}
\toprule
\textbf{Translation} & the language \\
\textbf{Root} & \textarabic{ل-غ-و} (l-ɣ-w) \\
\textbf{Pattern} & \textarabic{فُعْلَة} (fuʕlatun) \\
\textbf{Grammar} & Noun, feminine, genitive case (in iḍāfa construction), definite \\
\midrule
\textbf{Examples} & \makecell[l]{\parbox{9.5cm}{
1. \textarabic{أَتَكَلَّمُ ثَلَاثَ لُغَاتٍ} - I speak three languages [ʔatakallamu θalaːθa luɣaːtin]\\
2. \textarabic{اللُّغَةُ الْعَرَبِيَّةُ جَمِيلَةٌ} - The Arabic language is beautiful [al-luɣatu l-ʕarabijjatu ʤamiːlatun]\\
3. \textarabic{يَدْرُسُ لُغَةَ الْقُرْآنِ} - He studies the language of the Quran [jadrusu luɣata l-qurʔaːni]
}} \\
\midrule
\textbf{Synonyms} & \textarabic{لِسَان} (tongue/language), \textarabic{كَلَام} (speech) \\
\textbf{Etymology} & From root meaning "to speak idly, to be eloquent"; related to Hebrew \texthebrew{לשון} /laʃon/ \\
\bottomrule
\end{tabular}

\subsubsection*{Full Declension of \textarabic{اللُّغَة}}

\begin{tabular}{|c|c|c|c|}
\hline
\textbf{Number} & \textbf{Case} & \textbf{Indefinite} & \textbf{Definite} \\
\hline
\multirow{3}{*}{Singular (\textarabic{مُفْرَد})}
& Nominative & \textarabic{لُغَةٌ} /luɣatun/ & \textarabic{اللُّغَةُ} /al-luɣatu/ \\
& Accusative & \textarabic{لُغَةً} /luɣatan/ & \textarabic{اللُّغَةَ} /al-luɣata/ \\
& Genitive & \textarabic{لُغَةٍ} /luɣatin/ & \textarabic{اللُّغَةِ} /al-luɣati/ \\
\hline
\multirow{3}{*}{Dual (\textarabic{مُثَنَّى})}
& Nominative & \textarabic{لُغَتَانِ} /luɣataːni/ & \textarabic{اللُّغَتَانِ} /al-luɣataːni/ \\
& Acc/Gen & \textarabic{لُغَتَيْنِ} /luɣatajni/ & \textarabic{اللُّغَتَيْنِ} /al-luɣatajni/ \\
\hline
\multirow{3}{*}{Plural (\textarabic{جَمْع})}
& Nominative & \textarabic{لُغَاتٌ} /luɣaːtun/ & \textarabic{اللُّغَاتُ} /al-luɣaːtu/ \\
& Accusative & \textarabic{لُغَاتٍ} /luɣaːtin/ & \textarabic{اللُّغَاتِ} /al-luɣaːti/ \\
& Genitive & \textarabic{لُغَاتٍ} /luɣaːtin/ & \textarabic{اللُّغَاتِ} /al-luɣaːti/ \\
\hline
\end{tabular}

\subsection{\textarabic{الْعَرَبِيَّةِ} – [al-ʕarabijjati]}

\begin{tabular}{p{3cm}p{10cm}}
\toprule
\textbf{Translation} & the Arabic / the Arabian \\
\textbf{Root} & \textarabic{ع-ر-ب} (ʕ-r-b) \\
\textbf{Pattern} & \textarabic{فَعَلِيَّة} (faʕaliːja) - nisba adjective with feminine ending \\
\textbf{Grammar} & Adjective, feminine, genitive case (agreeing with \textarabic{اللُّغَةِ}), definite \\
\midrule
\textbf{Formation} & Nisba adjective: \textarabic{عَرَب} + \textarabic{ـِيّ} (nisba suffix) + \textarabic{ـَة} (feminine) \\
\midrule
\textbf{Examples} & \makecell[l]{\parbox{9.5cm}{
1. \textarabic{الثَّقَافَةُ الْعَرَبِيَّةُ} - Arab culture [aθ-θaqaːfatu l-ʕarabijjatu]\\
2. \textarabic{الْجَزِيرَةُ الْعَرَبِيَّةُ} - The Arabian Peninsula [al-ʤaziːratu l-ʕarabijjatu]\\
3. \textarabic{الْأَدَبُ الْعَرَبِيُّ} - Arabic literature (masc.) [al-ʔadabu l-ʕarabijju]
}} \\
\midrule
\textbf{Related Forms} & \textarabic{عَرَبٌ} (Arabs), \textarabic{عَرَبِيٌّ} (an Arab, masc.), \textarabic{عَرَبِيَّة} (an Arab woman) \\
\textbf{Etymology} & From \textarabic{عَرَب} meaning "Arabs, Bedouins, clear expression" \\
\bottomrule
\end{tabular}

\subsubsection*{Full Declension of \textarabic{عَرَبِيّ}}

\begin{tabular}{|c|c|c|c|c|}
\hline
\textbf{Gender} & \textbf{Case} & \textbf{Indefinite} & \textbf{Definite} \\
\hline
\multirow{3}{*}{Masculine}
& Nominative & \textarabic{عَرَبِيٌّ} /ʕarabijjun/ & \textarabic{الْعَرَبِيُّ} /al-ʕarabijju/ \\
& Accusative & \textarabic{عَرَبِيًّا} /ʕarabijjan/ & \textarabic{الْعَرَبِيَّ} /al-ʕarabijja/ \\
& Genitive & \textarabic{عَرَبِيٍّ} /ʕarabijjin/ & \textarabic{الْعَرَبِيِّ} /al-ʕarabijji/ \\
\hline
\multirow{3}{*}{Feminine}
& Nominative & \textarabic{عَرَبِيَّةٌ} /ʕarabijjatun/ & \textarabic{الْعَرَبِيَّةُ} /al-ʕarabijjatu/ \\
& Accusative & \textarabic{عَرَبِيَّةً} /ʕarabijjatan/ & \textarabic{الْعَرَبِيَّةَ} /al-ʕarabijjata/ \\
& Genitive & \textarabic{عَرَبِيَّةٍ} /ʕarabijjatin/ & \textarabic{الْعَرَبِيَّةِ} /al-ʕarabijjati/ \\
\hline
\end{tabular}

% ======================== Phrase Analysis ========================
\section{Phrase Analysis}
\begin{tcolorbox}[colback=boxcolor,colframe=headercolor,breakable]
\textbf{Grammatical Structure:}\\
Interrogative particle + present verb (impersonal) + prepositional phrase with pronoun + verbal noun (subject) + demonstrative pronoun (non-standard) + noun in genitive (iḍāfa) + adjective in genitive \\[0.5em]

\textbf{Key Grammar Points:}
\begin{itemize}
\item \textbf{Interrogative \textarabic{لِمَ}}: Compound of \textarabic{لِ} (for) + \textarabic{مَا} (what), with elision of alif
\item \textbf{Impersonal construction}: \textarabic{يَجِبُ عَلَيَّ} expresses necessity/obligation on the speaker
\item \textbf{Verbal noun as subject}: \textarabic{تَعَلُّمُ} (learning) functions as the grammatical subject
\item \textbf{Iḍāfa (construct state)}: \textarabic{تَعَلُّمُ هٰذِهِ اللُّغَةِ} - "learning of this language"
\item \textbf{Adjective agreement}: \textarabic{الْعَرَبِيَّةِ} agrees with \textarabic{اللُّغَةِ} in gender, case, and definiteness
\item \textbf{Non-standard demonstrative}: \textarabic{ذِي} is dialectal; standard form is \textarabic{هٰذِهِ}
\item \textbf{Full parsing}: \\
  - \textarabic{لِمَ}: interrogative particle (for what reason?)\\
  - \textarabic{يَجِبُ}: imperfect verb, 3rd masc. sing., indicative\\
  - \textarabic{عَلَيَّ}: preposition + 1st person pronoun\\
  - \textarabic{تَعَلُّمُ}: verbal noun, nominative (subject)\\
  - \textarabic{ذِي/هٰذِهِ}: demonstrative, genitive\\
  - \textarabic{اللُّغَةِ}: noun, genitive (in iḍāfa)\\
  - \textarabic{الْعَرَبِيَّةِ}: adjective, genitive
\end{itemize}
\end{tcolorbox}

\begin{tcolorbox}[colback=white,colframe=accentcolor,breakable,title=\textbf{Syntactic Analysis}]
\textbf{Sentence Type:} Interrogative (question)\\
\textbf{Question Type:} Asking for reason/justification (why?)\\
\textbf{Verb Pattern:} Impersonal verbal sentence with delayed subject\\
\textbf{Word Order:} VSO modified to: Interrogative + V + Prepositional Phrase + S\\
\textbf{Definiteness Chain:} \textarabic{هٰذِهِ اللُّغَةِ الْعَرَبِيَّةِ} - all definite through agreement
\end{tcolorbox}

% ======================== Similar Phrases for Practice ========================
\section{Similar Phrases for Practice}

\begin{enumerate}
\item \textarabic{لِمَ يَجِبُ عَلَيَّ تَعَلُّمُ الرِّيَاضِيَّاتِ؟}\\
Why should I learn mathematics? [lima jaʤibu ʕalajja taʕallumu r-rijaːdˤijjaːti]

\item \textarabic{لِمَ يَجِبُ عَلَيْكَ دِرَاسَةُ التَّارِيخِ؟}\\
Why should you study history? [lima jaʤibu ʕalajka diraːsatu t-taːriːxi]

\item \textarabic{لِمَاذَا يَنْبَغِي لَنَا أَنْ نَتَعَلَّمَ اللُّغَاتِ الْأَجْنَبِيَّةَ؟}\\
Why should we learn foreign languages? [limaːðaː janbaɣiː lanaː ʔan nataʕallama l-luɣaːti l-ʔaʤnabijjata]

\item \textarabic{لِمَ عَلَيَّ أَنْ أَدْرُسَ هٰذَا الْمَوْضُوعَ؟}\\
Why do I have to study this subject? [lima ʕalajja ʔan ʔadrusa haːðaː l-mawdˤuːʕa]

\item \textarabic{لِمَ يَلْزَمُنِي تَعَلُّمُ الْبَرْمَجَةِ؟}\\
Why is it necessary for me to learn programming? [lima jalzamuniː taʕallumu l-barʤamati]

\item \textarabic{لِأَيِّ سَبَبٍ يَجِبُ أَنْ أَتَعَلَّمَ الْفَلْسَفَةَ؟}\\
For what reason should I learn philosophy? [liʔajji sababin jaʤibu ʔan ʔataʕallama l-falsafata]
\end{enumerate}

% ======================== Levantine Dialect Version ========================
\section{Levantine (Shaami) Arabic Dialect}

\begin{tcolorbox}[colback=white,colframe=dialectcolor,title=\textbf{Levantine Version},breakable]
\textarabic{لَيْشْ لَازِم أَتْعَلَّمْ هَيْ اللُّغَة الْعَرَبِيَّة؟}\\
\textbf{Phonetic:} [leːʃ laːzim ʔatʕallam haj il-luɣa l-ʕarabijje]\\
\textbf{Translation:} Why do I have to learn this Arabic language?
\end{tcolorbox}

\textbf{Key Dialectal Changes:}
\begin{itemize}
\item \textarabic{لِمَ} → \textarabic{لَيْشْ} (leːʃ) – common Levantine "why" (from \textarabic{لِأَيِّ شَيْءٍ})
\item \textarabic{يَجِبُ عَلَيَّ} → \textarabic{لَازِم} (laːzim) – "it's necessary/I must"
\item \textarabic{تَعَلُّمُ} → \textarabic{أَتْعَلَّمْ} (ʔatʕallam) – first person present instead of verbal noun
\item \textarabic{هٰذِهِ} → \textarabic{هَيْ} (haj) – dialectal demonstrative "this" (feminine)
\item \textbf{Word order change}: Levantine uses interrogative + verb + subject, more SVO-like
\item Final case vowels are dropped in dialect (no tanwīn or iʕrāb)
\item \textarabic{الْ} becomes \textarabic{إِلْ} in connected speech
\end{itemize}

\subsection{Alternative Dialectal Variants}

\begin{tabular}{|l|l|l|}
\hline
\textbf{Dialect} & \textbf{Phrase} & \textbf{Phonetic} \\
\hline
Egyptian & \textarabic{لِيهْ لَازِمْ أَتْعَلِّمْ الْعَرَبِيّ؟} & [leː laːzim ʔatʕallim il-ʕarabiː] \\
\hline
Gulf & \textarabic{لَيْشْ لَازِمْ أَتْعَلَّمْ الْعَرَبِيّ؟} & [leːʃ laːzim ʔatʕallam il-ʕarabiː] \\
\hline
Moroccan & \textarabic{عْلَاشْ خَصْنِي نْتْعَلَّمْ الْعَرْبِيَّة؟} & [ʕlaːʃ xasˤniː ntʕallem l-ʕarbijja] \\
\hline
Iraqi & \textarabic{لَيْشْ لَازِمْ أَتْعَلَّمْ عَرَبِيّ؟} & [leːʃ laːzim ʔatʕallam ʕarabiː] \\
\hline
\end{tabular}

% ======================== Additional Learning Notes ========================
\section{Additional Learning Notes}

\begin{tcolorbox}[colback=boxcolor,colframe=accentcolor,title=\textbf{Cultural and Linguistic Context},breakable]
\textbf{Significance of Arabic:} Arabic is one of the six official languages of the United Nations and is spoken by over 400 million native speakers across 25+ countries. It's the liturgical language of Islam and has profoundly influenced languages from Spanish to Swahili.

\textbf{Classical vs. Modern Standard:} Modern Standard Arabic (MSA) is used in formal contexts, media, and literature. Classical Arabic is the language of the Quran and classical texts. Both share grammar but differ in vocabulary.

\textbf{Diglossia:} Arabic exhibits strong diglossia - formal MSA coexists with regional dialects that can be mutually unintelligible. Learning MSA provides access to all Arabic speakers, while dialects are used in daily conversation.

\textbf{Why Learn Arabic:}
\begin{itemize}
\item Access to rich literary, scientific, and philosophical heritage
\item Career opportunities in diplomacy, translation, international business
\item Understanding of Islamic texts and Middle Eastern culture
\item Cognitive benefits of learning a morphologically rich language
\item Gateway to understanding Semitic language family
\end{itemize}
\end{tcolorbox}

\subsection{Memory Tips}
\begin{enumerate}
\item \textbf{Question Formation:} Remember \textarabic{لِمَ} = \textarabic{لِ} (for) + \textarabic{مَا} (what) = "for what?" = "why?"
\item \textbf{Obligation Pattern:} \textarabic{يَجِبُ عَلَى} + pronoun is standard for "must/should"
\item \textbf{Form V Verbs:} Doubled middle radical (\textarabic{تَعَلَّمَ}) indicates reflexive/intensive action
\item \textbf{Nisba Adjectives:} Pattern \textarabic{فَعَلِيّ} creates relational adjectives (Arabic, Islamic, etc.)
\item \textbf{Iḍāfa Chain:} First noun loses nunation, following nouns in genitive
\item \textbf{Root System:} \textarabic{ع-ل-م} (knowledge) is one of the most productive roots in Arabic
\end{enumerate}

\subsection{Related Vocabulary Family}

\begin{multicols}{2}
\textbf{Education family:}\\
\textarabic{عِلْم} – knowledge [ʕilm]\\
\textarabic{مَعْرِفَة} – knowing [maʕrifa]\\
\textarabic{تَعْلِيم} – teaching [taʕliːm]\\
\textarabic{دِرَاسَة} – study [diraːsa]\\
\textarabic{طَالِب} – student [tˤaːlib]\\
\textarabic{مُعَلِّم} – teacher [muʕallim]\\
\textarabic{مَدْرَسَة} – school [madrasa]\\
\textarabic{جَامِعَة} – university [ʤaːmiʕa]\\

\textbf{Language family:}\\
\textarabic{لِسَان} – tongue/language [lisaːn]\\
\textarabic{كَلَام} – speech [kalaːm]\\
\textarabic{لَهْجَة} – dialect [lahʤa]\\
\textarabic{فُصْحَى} – classical/eloquent [fusˤħaː]\\
\textarabic{عَامِّيَّة} – colloquial [ʕaːmmijja]\\
\textarabic{نَحْو} – grammar [naħw]\\
\textarabic{صَرْف} – morphology [sˤarf]\\
\textarabic{مُفْرَدَات} – vocabulary [mufradaːt]\\

\textbf{Obligation family:}\\
\textarabic{وَاجِب} – obligatory [waːʤib]\\
\textarabic{لَازِم} – necessary [laːzim]\\
\textarabic{ضَرُورِيّ} – necessary [dˤaruːrij]\\
\textarabic{مَفْرُوض} – imposed [mafruːdˤ]\\
\textarabic{حَتْمِيّ} – inevitable [ħatmij]\\

\textbf{Question words:}\\
\textarabic{مَا/مَاذَا} – what [maː/maːðaː]\\
\textarabic{مَنْ} – who [man]\\
\textarabic{أَيْنَ} – where [ʔajna]\\
\textarabic{مَتَى} – when [mataː]\\
\textarabic{كَيْفَ} – how [kajfa]\\
\textarabic{كَمْ} – how many [kam]\\
\textarabic{أَيّ} – which [ʔajj]\\
\end{multicols}

\section{Grammatical Notes: The Construct State (Iḍāfa)}

\begin{tcolorbox}[colback=boxcolor,colframe=headercolor,breakable]
\textbf{Structure in this phrase:} \textarabic{تَعَلُّمُ هٰذِهِ اللُّغَةِ الْعَرَبِيَّةِ}

\textbf{Rules of Iḍāfa:}
\begin{enumerate}
\item First noun (muḍāf): \textarabic{تَعَلُّمُ} – loses nunation, cannot take \textarabic{الْ}
\item Second noun (muḍāf ilayhi): \textarabic{اللُّغَةِ} – must be in genitive case
\item Adjectives follow the entire construct and agree with the noun they modify
\item \textarabic{الْعَرَبِيَّةِ} modifies \textarabic{اللُّغَةِ}, not \textarabic{تَعَلُّمُ}
\item The demonstrative \textarabic{هٰذِهِ} can intervene in the iḍāfa as an "intruder"
\end{enumerate}

\textbf{Translation strategies:}
\begin{itemize}
\item Possessive: "learning's language" → "the learning of the language"
\item English prepositional phrase: "learning of this Arabic language"
\item Noun-noun compound: "Arabic language learning"
\end{itemize}
\end{tcolorbox}

\section{Answers to "Why Learn Arabic?"}

\begin{tcolorbox}[colback=white,colframe=dialectcolor,breakable]
\textbf{Possible answers in MSA:}

\begin{enumerate}
\item \textarabic{لِأَنَّ الْعَرَبِيَّةَ لُغَةٌ غَنِيَّةٌ وَجَمِيلَةٌ}\\
Because Arabic is a rich and beautiful language [liʔanna l-ʕarabij jata luɣatun ɣanijjatun wa ʤamiːlatun]

\item \textarabic{لِفَهْمِ الْقُرْآنِ وَالتُّرَاثِ الْإِسْلَامِيِّ}\\
To understand the Quran and Islamic heritage [lifahmi l-qurʔaːni wa t-turaːθi l-ʔislaːmijji]

\item \textarabic{لِتَوْسِيعِ آفَاقِي الثَّقَافِيَّةِ}\\
To broaden my cultural horizons [litawsiːʕi ʔaːfaːqiː θ-θaqaːfijjati]

\item \textarabic{لِأَنَّهَا تُفْتَحُ أَبْوَابَ الْفُرَصِ الْمِهْنِيَّةِ}\\
Because it opens doors of professional opportunities [liʔannahaː tuftaħu ʔabwaːba l-furasˤi l-mihnijjati]

\item \textarabic{لِلتَّوَاصُلِ مَعَ مَلَايِينِ النَّاطِقِينَ بِهَا}\\
To communicate with millions of its speakers [li t-tawaːsˤuli maʕa malaːjiːni n-naːtˤiqiːna bihaː]
\end{enumerate}
\end{tcolorbox}

\end{document}